\chapter{Données}
\label{chap:donnees}

Le projet a utilisé deux familles de données : un jeu \textbf{synthétique} produit par un
générateur physique, qui a servi à mettre au point la méthode, puis le domaine \textbf{réel} du
Grand Tunis, construit à partir de données ouvertes. Dans les deux cas, la vérité en altitude est
\emph{générée} : aucune mesure ne couvre le volume 3D. Nous parlons donc de \enquote{vérité
d'apprentissage} et réservons le mot \enquote{mesure} aux observations réelles.

\section{Le jeu synthétique ANEO}

\subsection{Génération}
Le générateur de l'écosystème ANEO produit, pour chaque instant, l'état de surface sur une grille
fine et le champ 3D correspondant, de façon cohérente avec la dynamique de la couche limite (vent
géostrophique, cisaillement, brise, relief). Le jeu \code{aneo\_dataset\_3d\_v2}, généré le
28 septembre 2026, est décrit au tableau~\ref{tab:synth}.

\begin{table}[H]
\centering
\caption{Caractéristiques du jeu synthétique \code{aneo\_dataset\_3d\_v2}.}
\label{tab:synth}
\begin{tabularx}{\linewidth}{@{}lX@{}}
\toprule
Élément & Valeur\\
\midrule
Durée & 100 jours, 1\,000 instants (environ 35\,Go)\\
Domaine & 50\,$\times$\,25\,km, relief jusqu'à 677\,m\\
Grille de surface & 1\,000\,$\times$\,500 points au pas de 50\,m\\
Nœuds 3D (cible) & 500\,$\times$\,250 au pas de 100\,m, 12 niveaux au-dessus du niveau de la mer : 25, 50, 75, 100, 150, 200, 300, 400, 500, 650, 800 et 1\,000\,m\\
Nœuds valides & 849\,184 au-dessus du relief (environ 1,5 million de nœuds de grille par instant)\\
Entrées de surface & vent $u$, $v$, $w$, température, densité, humidité relative, pression, et indicateur de mesure \code{mask\_i\_obs}\\
Cibles & $u$, $v$, $w$ et $\rho$ aux nœuds 3D\\
Découpage & par jour entier, 70 / 15 / 15\,\% (entraînement / validation / test)\\
Test & 166 instants : 120 synoptiques (heure pleine, multiple de 3\,h) et 46 d'épisodes rapides (rafales, front, brise), rapportés à part\\
\bottomrule
\end{tabularx}
\end{table}

Le contrôle qualité du générateur passe toutes ses vérifications (résidu de masse maximal
$6{,}4\cdot10^{-12}$, corrélations spatiales et temporelles, équilibre hydrostatique) sauf
l'équilibre des secteurs de direction (0,095 pour un seuil de 0,10), un défaut mineur non bloquant.

\subsection{Limites du jeu synthétique}
La couche de surface du générateur est quasi continue : environ 62\,000 points de mesure visibles
à 50\,\%, alors qu'un réseau ANEO réel compte un capteur tous les 500\,m à 1\,km. Le modèle entraîné
sur ce jeu voit donc beaucoup plus d'information qu'il n'en aura en production. Par ailleurs, la
position géographique du domaine synthétique est supposée : aucune mesure réelle ne peut y être
placée. Ces deux limites ont motivé le passage au domaine réel.

\subsection{Contrat de données}
Les fichiers de mesures ANEO suivent un contrat de données (critère DE1) : colonnes et types
attendus, plages physiques, règles de cohérence et indicateurs de qualité (valeurs hors plage,
figées, absentes, interpolées). Chaque fichier reçoit un verdict \enquote{conforme},
\enquote{quarantaine} ou \enquote{rejet} accompagné d'un rapport ; un fichier en quarantaine est
tout de même reconstruit, les stations signalées étant masquées.

\section{Le domaine réel du Grand Tunis}

\subsection{Grille et surface}
Le domaine réel couvre 68\,$\times$\,56\,km, de 36,395 à 36,900°\,N et de 10,04 à 10,80°\,E. Le
relief provient du modèle numérique de terrain Copernicus, et la rugosité $z_0$ de l'occupation du
sol ESA WorldCover \citep{zanaga2021}, sur une grille de surface au pas de 50\,m. Les nœuds 3D sont
au pas de 100\,m sur les douze niveaux du jeu synthétique : une grille de
12\,$\times$\,560\,$\times$\,680, dont 3\,469\,007 nœuds valides sur 4\,569\,600. La part du domaine
au-dessus du relief va de 26\,\% à 25\,m à 100\,\% à partir de 650\,m.

\subsection{Forçage ERA5 et vérité générée}
La vérité d'apprentissage est produite par le générateur \emph{forcé par ERA5} : pour chaque
instant, la grande échelle (vent géostrophique, cisaillement et rotation avec l'altitude, stabilité,
contraste de température terre--mer) est calculée à partir d'ERA5 sur les mailles qui couvrent le
domaine, puis le générateur y ajoute la brise de mer et l'effet du relief. La grande échelle est
recalée sur ERA5 maille par maille à la même hauteur au-dessus du sol, et $w$ est obtenu par
continuité. Deux portes de contrôle ont précédé la génération complète :
\begin{itemize}
  \item un jour d'hiver test : conservation de la masse exacte, moyenne du domaine égale à ERA5 à
        chaque niveau, erreur aux trois stations au plus égale à celle d'ERA5 ;
  \item le cycle diurne de la direction en été à Tunis-Carthage et Nabeul (brise de mer), validé
        visuellement avant la génération.
\end{itemize}

\textbf{Règle d'usage des stations} : les mesures des stations n'entrent jamais dans la vérité
d'apprentissage. Elles servent d'entrée au modèle et de référence de validation, ce qui garde la
validation sur le réel indépendante.

\subsection{Mesures réelles}
\begin{table}[H]
\centering
\caption{Sources de données réelles.}
\label{tab:sources}
\begin{tabularx}{\linewidth}{@{}lXl@{}}
\toprule
Source & Contenu & Usage\\
\midrule
METAR DTTA & Tunis-Carthage, vent à 10\,m, horaire & entrée, validation\\
SYNOP 60729, 60728 & Zaghouan et Nabeul (INM, via NOAA ISD), vent à 10\,m, triho\-raire & entrée, validation\\
IGRA TSM00060715 & Radiosondages de Tunis-Carthage & validation en altitude\\
ERA5, ERA5-Land & Réanalyse horaire, 0,25° et 0,1° & forçage, entrée du modèle\\
GFS & Prévision NOAA 0,25° (temps réel et archive AWS) & temps réel, OS7\\
GEFS & Ensemble de 31 membres & prévision probabiliste\\
Copernicus DEM, WorldCover & Relief et occupation du sol & grille, rugosité\\
\bottomrule
\end{tabularx}
\end{table}

Les séries ont été nettoyées : à Nabeul, l'anémomètre est resté à zéro en juin, juillet et
septembre 2021 (mois masqués), et les années 2023--2024, où la station est abritée, sont exclues.

\subsection{Découpage temporel}
Le découpage est temporel et fixé une fois pour toutes : le \textbf{test} couvre janvier--mars et
juillet--août 2022 (hiver et été), la \textbf{validation} janvier--février et juillet 2021, et
l'\textbf{entraînement} le reste de 2021 à 2025, avec une marge de sept jours autour des blocs de
validation et de test. Le cache du modèle compte 700 instants d'entraînement, 120 de validation et
180 de test (environ 49\,Mo par instant).

\subsection{Entrée de production}
Avec trois stations seulement dans le domaine, la couche de surface dense du jeu synthétique n'a
pas d'équivalent réel. L'entrée de production est donc formée de \textbf{pseudo-capteurs ERA5} à
10\,m tous les 2,5\,km, des \textbf{cellules des stations} réelles, des \textbf{sept scalaires de
forçage} de l'instant et du \textbf{profil ERA5} interpolé à chaque nœud. En temps réel et en
prévision, la prévision GFS remplace ERA5, avec les mêmes formules et les mêmes mailles.
